\section{Literature Review}\label{sec:2}

Recruitment technology, agentic software, and information retrieval have each advanced independently, yet deployed hiring tools rarely combine all three in a traceable architecture.
We organize prior work into five literature themes---recruitment systems, multi-agent coordination, semantic representation, ranking and fusion---then state the gaps those literatures leave open and how JobMatch is positioned relative to them.
Numeric results from our evaluation appear in \S\ref{sec:6}; metric definitions in \S\ref{sec:5}.

\subsection{Recruitment and job-matching systems}\label{sec:2-recruitment}

Enterprise hiring still depends heavily on manual screening: recruiters read resumes, compare them to written requirements, and decide who advances.
At scale this workflow is slow, inconsistent across reviewers, and weakly auditable when a candidate asks why they were excluded \todo{Citation needed: empirical studies of manual resume screening throughput and reviewer inconsistency in recruitment}.
Most production stacks therefore automate the first cut through keyword filters, structured application forms, and applicant tracking systems (ATS).

Keyword-based filtering treats a resume as a bag of terms matched against job requirements.
Lexical retrieval methods such as TF--IDF weighting and BM25 remain the workhorse for this pattern because they are fast, interpretable at the token level, and easy to deploy over large applicant pools~\cite{Manning2008,Robertson2009}.
Their limitation is semantic brittleness: synonyms, paraphrases, and transferable skills that are not spelled out literally can be missed, while keyword stuffing can inflate apparent fit.
Structured skill fields and Boolean search rules reduce noise but still encode fit as a static rule set rather than a graded comparison between two evolving profiles.

ATS platforms extend ingestion with workflow state---application status, interview stages, compliance logging---but they often treat matching as a side feature bolted onto document storage \todo{Citation needed: survey or industry analysis of applicant tracking system capabilities and matching modules}.
Candidates frequently experience opaque ranking: a profile passes a filter without knowing which requirement mattered, or disappears after a re-rank with no explanation.
Recruiters face the symmetric problem when shortlists shift without a visible link between job-text edits and ranking changes.
These workflow limitations motivate systems that separate \emph{profile ownership} from \emph{scoring}, and that return reasons a human can inspect before any hire/no-hire action is taken.

\subsection{AI agents and multi-agent systems}\label{sec:2-agents}

Agent-based software assigns bounded roles---perception, memory, planning, action---to modules that coordinate through messages or shared events rather than through one monolithic controller~\cite{Wooldridge2009}.
Three agent patterns are relevant to hiring.

\leadpara{Personal agents.}
Personal or client-side agents maintain a durable model of a user's goals and history, updating that model as new documents arrive \todo{Citation needed: personal software agents or digital assistants for career and document management}.
In hiring, the natural analogue is a representative that ingests resumes, tracks skill and experience revisions, and speaks for the candidate without exposing raw edits to the employer's datastore.

\leadpara{Task-specific agents.}
Task-specific agents encapsulate one domain workflow---parsing, validation, or recommendation---with explicit inputs and outputs.
Employer-side job-ingestion agents fit this pattern: they normalize postings, extract requirements, and version job records when descriptions change.
Prior agentic hiring proposals often emphasize autonomous interview bots or negotiation agents \todo{Citation needed: autonomous interview or negotiation agents in recruitment}.
JobMatch instead uses task-specific agents for \emph{profile maintenance}, not for autonomous selection.

\leadpara{Collaborative agents and human-in-the-loop support.}
Collaborative multi-agent systems decompose a problem so that each agent owns part of the state and a coordinator integrates read-only views for decision support~\cite{Wooldridge2009}.
Human-in-the-loop designs keep consequential actions---apply, reject, shortlist---with people while software proposes ranked options and explanations \todo{Citation needed: human-in-the-loop decision support for algorithmic hiring recommendations}.
This differs from fully automated ranking pipelines that update applicant status without an explicit user action on each match.

JobMatch adopts the collaborative pattern: Candidate and Employer agents own their respective profiles and vector stores; a read-only Matchmaking agent scores cross-side pairs and emits structured explanations (\S\ref{sec:3}).
An in-process event bus synchronizes profile updates and match-cache invalidation without cross-agent mutation of stored records (\S\ref{sec:4}).

\subsection{Semantic search and representation learning}\label{sec:2-semantic}

Job--candidate matching is, at core, a cross-document retrieval problem: given a resume (query), rank jobs (documents), or the reverse for employer-initiated search.
Semantic search replaces exact term overlap with learned representations that map text into a shared vector space~\cite{Manning2008}.

\leadpara{Embeddings.}
Dense encoders map resumes and job descriptions to fixed-dimensional vectors so that related texts lie nearby even when they share few tokens.
Sentence-level models such as Sentence-BERT train siamese encoders for semantic similarity and have become a default building block for short-text matching~\cite{Reimers2019}.
In JobMatch, snapshot embeddings are computed from canonical resume and job-description text (\S\ref{sec:4-params}).

\leadpara{Dense retrieval and semantic similarity.}
Dense retrieval ranks candidates or jobs by vector similarity---typically cosine---over the full corpus or over an approximate nearest-neighbor shortlist \todo{Citation needed: dense passage or document retrieval over learned embeddings}.
Compared with lexical baselines, dense retrieval captures paraphrase and topical relatedness but can underweight explicit skill tokens that recruiters treat as hard requirements.
Hybrid systems therefore retain both signals rather than replacing keywords entirely~\cite{Manning2008,Reimers2019}.

Much published job-matching research evaluates a single embedding model or similarity function on a static labeled set \todo{Citation needed: job--resume matching studies using neural or embedding-based retrieval}.
Fewer prototypes expose the same representation path to live profile edits on both market sides while reporting offline retrieval metrics on identical scoring code.

\subsection{Ranking and recommendation}\label{sec:2-ranking}

Once candidates and jobs are represented, the system must \emph{rank} pairs and \emph{evaluate} ranking quality---topics well developed in information retrieval and recommender-systems research.

\leadpara{Ranking metrics.}
Standard offline metrics include precision and recall at cutoff $K$, mean reciprocal rank (MRR), normalized discounted cumulative gain (nDCG@K), and mean average precision (MAP)~\cite{Manning2008}.
Precision@K measures shortlist purity; recall@K measures coverage of labeled relevant items; nDCG@K rewards placing higher-grade relevance labels earlier; MRR emphasizes the rank of the first relevant result.
Graded nDCG is appropriate when relevance is partial (e.g., a job is a plausible but not ideal fit).
We adopt these metrics at $K{=}5$ on a labeled demo corpus (\S\ref{sec:5}).

\leadpara{Hybrid retrieval.}
Hybrid rankers combine lexical scores, dense semantic similarity, and structured features such as skill overlap or experience gaps~\cite{Manning2008,Robertson2009,Reimers2019}.
A common pattern is a weighted sum of semantic and skill channels (multimodal blending) or a post-hoc constraint filter on top of a semantic shortlist.
JobMatch extends this with a multi-component composite score---semantic, skills, title overlap, experience, compensation, and remote preference---so that each channel remains visible in portal explanations (\S\ref{sec:5}).

\leadpara{Fusion methods.}
When multiple rankers produce different orderings, fusion methods merge lists without retraining the underlying encoders.
Reciprocal rank fusion (RRF) combines ranked outputs using reciprocal rank weights and remains stable when individual retrievers make complementary errors~\cite{Cormack2009}.
Learned fusion and hierarchical blending fit scalar weights or classifiers over component scores \todo{Citation needed: learned rank fusion or stacking for multi-signal retrieval}.
JobMatch implements fixed-weight composite scoring for portal defaults, plus optional RRF, learned logistic fusion, Platt calibration~\cite{Platt1999}, and cross-encoder reranking for offline comparison (\S\ref{sec:4}, \S\ref{sec:6}).

\subsection{Research gaps}\label{sec:2-gaps}

Across the recruitment, agent, and retrieval literatures, several gaps recur in job-matching demos and prototypes.

\leadpara{Algorithm-only focus.}
Many systems optimize a matching formula---embedding similarity, skill Jaccard, or a learned re-ranker---while treating ingestion, profile versioning, and user-facing workflow as out of scope \todo{Citation needed: job-matching prototypes that report retrieval metrics without deployable profile or portal layers}.
Evaluation then proceeds on frozen corpora disconnected from the code path users would actually trigger when they edit a resume or repost a job.

\leadpara{Single-sided or monolithic control.}
Candidate-side recommender demos are more common than symmetric platforms where \emph{both} sides maintain agents with separate ownership \todo{Citation needed: two-sided or employer--candidate job-matching platforms with separate user roles}.
Monolithic services that store and score all records in one module obscure accountability: users cannot tell which component changed a rank or which side's edit caused invalidation.

\leadpara{Limited explainability and evaluation depth.}
Ranking scores alone do not satisfy recruiter or candidate trust; explainable ML methods attribute predictions to input features~\cite{Ribeiro2016}.
In hiring, explanations should reference skills, experience gaps, and constraint channels that align with the score breakdown.
Yet many job-matching papers report nDCG or accuracy without explanation-quality audits, calibration checks~\cite{Platt1999}, or fairness sanity tests on controlled profile variants \todo{Citation needed: explainability or fairness evaluation in job-matching or hiring recommender systems}.
Demo corpora are often small and manually labeled, which is appropriate for reproducible research but rarely accompanied by regression gates tied to production scoring code.

\subsection{Positioning of JobMatch}\label{sec:2-positioning}

JobMatch addresses these gaps as an integrated research prototype rather than as a standalone retrieval formula.

\leadpara{Multi-agent framing.}
The system decomposes hiring into three bounded agents---Candidate, Employer, and Matchmaking---coordinated by typed events (\S\ref{sec:3}--\S\ref{sec:4}).
Profile mutation stays with the owning agent; matching reads versioned snapshots only.
This design makes ownership, cache invalidation, and audit trails first-class, not emergent properties of a single service.

\leadpara{Candidate-side and employer-side portals.}
Separate React portals enforce role separation: candidates manage resumes, saved jobs, and applications; employers manage postings, reverse candidate search, and applicant lists; administrators inspect agent status, run evaluation hooks, and review fairness summaries (\S\ref{sec:4-frontend}).
Consequential actions remain explicit user clicks---save, apply, dismiss---consistent with human-in-the-loop decision support.

\leadpara{Matchmaking agent and explainable composite ranking.}
The Matchmaking agent combines semantic, skill, and constraint channels into a fixed composite score with structured \texttt{MatchExplanation} payloads (\S\ref{sec:5}).
Rule-based explainers are evaluated offline for faithfulness and specificity (\S\ref{sec:5-explainability}); optional calibration and feedback boosts exist but are disabled in default portal settings.

\leadpara{Empirical evaluation on shared code paths.}
Offline benchmarks call the same scoring core as the live gateway: lexical and dense baselines, multimodal and composite strategies, RRF and learned fusion, ablations, hard-negative sanity checks, synthetic fairness counterfactuals, and a continuous-integration regression gate on nDCG@5 (\S\ref{sec:5}--\S\ref{sec:6}).
We do not claim state-of-the-art hiring accuracy on a large production corpus; the contribution is traceable multi-agent architecture, assistive explainable ranking, and reproducible measurement on a controlled labeled set tied to the deployed prototype.
